\section{实验设置}
\label{sec:experiments-zh}

\subsection{语料与划分}

程序化语料包含 24,000 个具有显式保存归属的片段，其中训练集 20,000 个、验证集 2,000 个、测试集 2,000 个。生成使用随机种子 42，片段包含 4--16 小节、2--8 个声部、六种乐器标签、七种节奏轮廓和七种姿态标签。划分摘要记录 \texttt{smoke=false}。CPU 冒烟样本单独保存，不进入正式指标。

\subsection{训练方案}

所有正式神经配置均使用六层 Transformer，隐藏维度为 384，注意力头数为 6，前馈维度为 1,536，dropout 为 0.1，上下文长度为 256。优化器为 AdamW，学习率 $3\times10^{-4}$，权重衰减 0.01，物理批量大小 16，梯度裁剪范数 1.0；CUDA 环境下启用 fp16 自动混合精度，Windows 数据加载线程数为 0。训练上限为 60 轮，早停耐心值为 10。若 CUDA 显存不足，运行器依次尝试批量大小 8、梯度累积 2 步，以及批量大小 4、梯度累积 4 步，且不缩减语料。

每轮训练写入可恢复的最后检查点，其中包含模型、优化器、缩放器、随机数状态、轮次历史以及语料和配置哈希；最佳验证检查点另行保存。仅当检查点及已完成训练摘要与当前配置和处理后语料哈希一致时，才允许正式评估。

\subsection{硬件与软件门控}

正式实验运行于 Windows 11、Python 3.11.9、PyTorch 2.5.1+cu121 和 CUDA 环境，GPU 为具有 16GB 显存的 NVIDIA GeForce RTX 4060 Ti。所有已完成训练的实测进程内存峰值为 4.637GiB，CUDA 已分配显存峰值为 0.667GiB。实验未发生需要调整批量大小的显存不足问题。

\subsection{比较配置}

完整运行器评估了 13 个具名配置：一个独立程序化规则参照、一个使用单候选解码的共享神经生成器、同一生成器的四候选约束引导解码、两个无约束对照、四个单条件移除模型和四个聚焦条件模型。所有神经条件消融均使用相同的已保存语料与划分。规则参照为每个测试条件生成新的确定性实现，不读取保存的目标事件列表。

拟议模型与普通模型共享一个已训练条件生成器，但二者比较同时把候选预算从 1 改为 4，并启用符号候选评分。消融模型分别训练，隐藏指定的条件前缀字段，并采用与普通模型相同的单候选解码预算；评估目标保持不变。矩阵包含 12 个神经检查点路径，但只有 11 次不同的神经优化运行，因为普通模型与拟议模型参数完全相同。所有神经配置均达到 60 轮上限，记录的最佳轮次介于 56 至 60。

\subsection{评估与导出}

教师强制标记指标枚举 2,000 个测试样本中的全部乐谱主体窗口。每个配置的生成指标使用 2,000 个测试条件。每个已评估配置尝试导出 20 个 MusicXML，最终专家包只接收结构可解析且与请求小节数和声部数一致的文件。13 行指标、全部所需检查点以及 260 条尝试导出记录均通过运行器最终门控。旧版 v2 检查点、冒烟输出和历史多种子开发实验不进入本文结果。
